% GENERATED FILE. Edit canonical lessons, then run npm run generate:books.
% canonical-source-hash: fnv1a64:3cb33900953b847b
% canonical-lessons: TA-C324-ancali, TA-C324-irutimariyatai, TA-C324-ninaivunal, TA-C324-eriyuttu, TA-C324-ninaivental

\chapter{Remembering the Dead}
\label{ch:ta-remembering-the-dead}
\hlchaptermodality{324}

\begin{chapteropening}
\textbf{By the end of this chapter:} \emph{I can say how people remember someone who died.}

Complete the last lesson of chapter 324: I can say how people remember someone who died.
\end{chapteropening}

\section[\texorpdfstring{\ta{அஞ்சலி} (añcali)}{அஞ்சலி (añcali)}]{\ta{அஞ்சலி} (añcali) — a tribute; homage}

\label{lesson:TA-C324-ancali}



\begin{quote}
Before the new one: say the Tamil for to bury, then the Tamil for a grave.
\end{quote}

\subsection*{You'll want to know: \ta{அஞ்சலி}}
\textbf{\ta{அஞ்சலி}} — \emph{añcali} — \textquotedblleft{}a tribute; homage\textquotedblright{}.

\begin{quote}
\ta{மக்கள்} \ta{அவருக்கு} \ta{அஞ்சலி} \ta{செலுத்தினார்கள்}.
\end{quote}

\begin{quote}
\emph{makkaḷ avarukku añcali celuttiṉārkaḷ.}
\end{quote}

\begin{quote}
The people paid tribute to him.
\end{quote}

\ta{அஞ்சலி} \ta{செலுத்து}, to pay homage; the person honoured takes -\ta{க்கு}.

\begin{cousinweb}
\textbf{\ta{அஞ்சலி}} (\emph{añcali}) is borrowed from Sanskrit \emph{añjali}, \textquotedblleft{}the hollowed hands placed together\textquotedblright{} (CDIAL 171), the gesture of homage.
\end{cousinweb}

\begin{grammarlens}[title={What you've built}]
One more word for telling what happened.
\end{grammarlens}

\subsection*{Your turn}
\begin{itemize}
\raggedright
  \item \emph{Say it:} \emph{añcali}
  \item \emph{Say it:} \emph{añcali}, once more
  \item \emph{Say it:} \emph{makkaḷ avarukku añcali celuttiṉārkaḷ.} — The people paid tribute to him.
  \item \emph{Recall:} say the Tamil for to bury, then the Tamil for a grave, then say \emph{añcali} again
\end{itemize}

\begin{tcolorbox}[breakable,colback=teal!4,colframe=teal!35!black,title={Before you move on}]
What does \ta{அஞ்சலி} mean? (\textquotedblleft{}A tribute; homage\textquotedblright{}.) Say it once more.
\end{tcolorbox}
\section[\texorpdfstring{\ta{இறுதிமரியாதை} (iṟutimariyātai)}{இறுதிமரியாதை (iṟutimariyātai)}]{\ta{இறுதிமரியாதை} (iṟutimariyātai) — last respects}

\label{lesson:TA-C324-irutimariyatai}



\begin{quote}
Before the new one: say the Tamil for a grave, then the Tamil for a tribute.
\end{quote}

\subsection*{You'll want to know: \ta{இறுதிமரியாதை}}
\textbf{\ta{இறுதிமரியாதை}} — \emph{iṟutimariyātai} — \textquotedblleft{}last respects\textquotedblright{}.

\begin{quote}
\ta{ஊர்} \ta{மக்கள்} \ta{அவருக்கு} \ta{இறுதிமரியாதை} \ta{செலுத்தினார்கள்}.
\end{quote}

\begin{quote}
\emph{ūr makkaḷ avarukku iṟutimariyātai celuttiṉārkaḷ.}
\end{quote}

\begin{quote}
The villagers paid him their last respects.
\end{quote}

\ta{இறுதி}, last + \ta{மரியாதை}, respect.

\begin{cousinweb}
\textbf{\ta{இறுதிமரியாதை}} (\emph{iṟutimariyātai}) joins native \ta{இறுதி}, \textquotedblleft{}end\textquotedblright{} (DEDR 514), with \ta{மரியாதை}, \textquotedblleft{}respect\textquotedblright{}, borrowed from Sanskrit \emph{maryādā}, \textquotedblleft{}boundary, propriety\textquotedblright{} (CDIAL 9895).
\end{cousinweb}

\begin{grammarlens}[title={What you've built}]
One more word for telling what happened.
\end{grammarlens}

\subsection*{Your turn}
\begin{itemize}
\raggedright
  \item \emph{Say it:} \emph{iṟutimariyātai}
  \item \emph{Say it:} \emph{iṟutimariyātai}, once more
  \item \emph{Say it:} \emph{ūr makkaḷ avarukku iṟutimariyātai celuttiṉārkaḷ.} — The villagers paid him their last respects.
  \item \emph{Recall:} say the Tamil for a grave, then the Tamil for a tribute, then say \emph{iṟutimariyātai} again
\end{itemize}

\begin{tcolorbox}[breakable,colback=teal!4,colframe=teal!35!black,title={Before you move on}]
What does \ta{இறுதிமரியாதை} mean? (\textquotedblleft{}Last respects\textquotedblright{}.) Say it once more.
\end{tcolorbox}
\section[\texorpdfstring{\ta{நினைவுநாள்} (niṉaivunāḷ)}{நினைவுநாள் (niṉaivunāḷ)}]{\ta{நினைவுநாள்} (niṉaivunāḷ) — a memorial day}

\label{lesson:TA-C324-ninaivunal}



\begin{quote}
Before the new one: say the Tamil for a tribute, then the Tamil for last respects.
\end{quote}

\subsection*{You'll want to know: \ta{நினைவுநாள்}}
\textbf{\ta{நினைவுநாள்}} — \emph{niṉaivunāḷ} — \textquotedblleft{}a memorial day\textquotedblright{}.

\begin{quote}
\ta{இன்று} \ta{என்} \ta{தாத்தாவின்} \ta{நினைவுநாள்}.
\end{quote}

\begin{quote}
\emph{iṉṟu eṉ tāttāviṉ niṉaivunāḷ.}
\end{quote}

\begin{quote}
Today is my grandfather's memorial day.
\end{quote}

The yearly day of a death, kept in memory.

\begin{cousinweb}
\textbf{\ta{நினைவுநாள்}} (\emph{niṉaivunāḷ}) is a Tamil compound built from \ta{நினைவு}, \textquotedblleft{}remembrance\textquotedblright{} (DEDR 3683), and \ta{நாள்}, \textquotedblleft{}day\textquotedblright{} (DEDR 3656).
\end{cousinweb}

\begin{grammarlens}[title={What you've built}]
One more word for telling what happened.
\end{grammarlens}

\subsection*{Your turn}
\begin{itemize}
\raggedright
  \item \emph{Say it:} \emph{niṉaivunāḷ}
  \item \emph{Say it:} \emph{niṉaivunāḷ}, once more
  \item \emph{Say it:} \emph{iṉṟu eṉ tāttāviṉ niṉaivunāḷ.} — Today is my grandfather's memorial day.
  \item \emph{Recall:} say the Tamil for a tribute, then the Tamil for last respects, then say \emph{niṉaivunāḷ} again
\end{itemize}

\begin{tcolorbox}[breakable,colback=teal!4,colframe=teal!35!black,title={Before you move on}]
What does \ta{நினைவுநாள்} mean? (\textquotedblleft{}A memorial day\textquotedblright{}.) Say it once more.
\end{tcolorbox}
\section[\texorpdfstring{\ta{எரியூட்டு} (eriyūṭṭu)}{எரியூட்டு (eriyūṭṭu)}]{\ta{எரியூட்டு} (eriyūṭṭu) — to cremate}

\label{lesson:TA-C324-eriyuttu}



\begin{quote}
Before the new one: say the Tamil for last respects, then the Tamil for a memorial day.
\end{quote}

\subsection*{You'll want to know: \ta{எரியூட்டு}}
\textbf{\ta{எரியூட்டு}} — \emph{eriyūṭṭu} — \textquotedblleft{}to cremate\textquotedblright{}.

\begin{quote}
\ta{அவருடைய} \ta{உடலை} \ta{மாலையில்} \ta{எரியூட்டினார்கள்}.
\end{quote}

\begin{quote}
\emph{avaruṭaiya uṭalai mālaiyil eriyūṭṭiṉārkaḷ.}
\end{quote}

\begin{quote}
They cremated his body in the evening.
\end{quote}

Past (I …): \ta{எரியூட்டினேன்}, \emph{eriyūṭṭiṉēṉ}.

Said of a cremation; the body takes -\ta{ஐ}.

\begin{cousinweb}
\textbf{\ta{எரியூட்டு}} (\emph{eriyūṭṭu}) is a Tamil compound built from \ta{எரி}, \textquotedblleft{}fire, flame\textquotedblright{} (DEDR 811), and \ta{ஊட்டு}, \textquotedblleft{}to feed\textquotedblright{} (DEDR 600).
\end{cousinweb}

\begin{grammarlens}[title={What you've built}]
One more word for telling what happened.
\end{grammarlens}

\subsection*{Your turn}
\begin{itemize}
\raggedright
  \item \emph{Say it:} \emph{eriyūṭṭu}
  \item \emph{Say it:} \emph{eriyūṭṭu}, once more
  \item \emph{Say it:} \emph{avaruṭaiya uṭalai mālaiyil eriyūṭṭiṉārkaḷ.} — They cremated his body in the evening.
  \item \emph{Recall:} say the Tamil for last respects, then the Tamil for a memorial day, then say \emph{eriyūṭṭu} again
\end{itemize}

\begin{tcolorbox}[breakable,colback=teal!4,colframe=teal!35!black,title={Before you move on}]
What does \ta{எரியூட்டு} mean? (\textquotedblleft{}To cremate\textquotedblright{}.) Say it once more.
\end{tcolorbox}
\section[\texorpdfstring{\ta{நினைவேந்தல்} (niṉaivēntal)}{நினைவேந்தல் (niṉaivēntal)}]{\ta{நினைவேந்தல்} (niṉaivēntal) — a memorial meeting}

\label{lesson:TA-C324-ninaivental}



\begin{quote}
Before the new one: say the Tamil for a memorial day, then the Tamil for to cremate.
\end{quote}

\subsection*{You'll want to know: \ta{நினைவேந்தல்}}
\textbf{\ta{நினைவேந்தல்}} — \emph{niṉaivēntal} — \textquotedblleft{}a memorial meeting\textquotedblright{}.

\begin{quote}
\ta{பள்ளியில்} \ta{ஒரு} \ta{நினைவேந்தல்} \ta{கூட்டம்} \ta{நடந்தது}.
\end{quote}

\begin{quote}
\emph{paḷḷiyil oru niṉaivēntal kūṭṭam naṭantatu.}
\end{quote}

\begin{quote}
A memorial meeting was held at the school.
\end{quote}

\ta{நினைவு}, memory + \ta{ஏந்தல்}, holding up.

\begin{cousinweb}
\textbf{\ta{நினைவேந்தல்}} (\emph{niṉaivēntal}) is a Tamil compound built from \ta{நினைவு}, \textquotedblleft{}remembrance\textquotedblright{} (DEDR 3683), and \ta{ஏந்தல்}, \textquotedblleft{}holding up\textquotedblright{}, from \ta{ஏந்து}, \textquotedblleft{}to hold up in the hands\textquotedblright{} (DEDR 894).
\end{cousinweb}

\begin{grammarlens}[title={What you've built}]
One more word for telling what happened.
\end{grammarlens}

\subsection*{Your turn}
\begin{itemize}
\raggedright
  \item \emph{Say it:} \emph{niṉaivēntal}
  \item \emph{Say it:} \emph{niṉaivēntal}, once more
  \item \emph{Say it:} \emph{paḷḷiyil oru niṉaivēntal kūṭṭam naṭantatu.} — A memorial meeting was held at the school.
  \item \emph{Recall:} say the Tamil for a memorial day, then the Tamil for to cremate, then say \emph{niṉaivēntal} again
\end{itemize}

\begin{tcolorbox}[breakable,colback=teal!4,colframe=teal!35!black,title={Before you move on}]
What does \ta{நினைவேந்தல்} mean? (\textquotedblleft{}A memorial meeting\textquotedblright{}.) Say it once more.
\end{tcolorbox}
